\begin{theorem}[Regular closures and local parent conditions]%
    \label{thm:peps_regular_closure_parent_constraint}
    \lean{TNLean.PEPS.IsGInjective.commute_of_cornerParent_annihilates_torusGClosure,
        TNLean.PEPS.IsGInjective.regionParent_annihilates_torusGClosure_of_isSimplyConnected,
        TNLean.PEPS.IsGInjective.regionParentHamiltonian_annihilates_torusGClosure,
        TNLean.PEPS.IsGInjective.span_torusGClosure_commuting_le_parentKernel,
        TNLean.PEPS.IsGInjective.card_commutingPairConjugacyClass_le_finrank_parentKernel}
    \leanok
    \uses{thm:peps_regular_torus_parent_flatness,
        thm:peps_regular_region_parent_flatness,
        def:peps_region_parent_interaction, thm:peps_pair_centralizer_count}
    Let both torus periods be at least three, and let the native
    tensor be regular $G$-injective. The original parent condition
    on the plaquette at the intersection of the two seams forces
    the labels $g,h$ of a closure vector to commute. Conversely,
    every commuting closure satisfies the original parent condition
    on every connected region whose actual closed-cell realization
    is simply connected. These regions may meet either seam.

    For any finite family of such regions, let $H$ be a sum of
    positive parent interactions with the original regional kernels.
    Then the span of the commuting closure vectors is contained in
    $\ker H$, and
    \[
        \dim\ker H\ \geq\
        \sum_{C\in\operatorname{Conj}(G)}
        \left|\operatorname{Conj}\bigl(Z_G(c_C)\bigr)\right|,
    \]
    where $c_C$ represents $C$. No $G$-isometry is required.
    This proves the inclusion and the resulting lower bound for
    the regular case of \cite[Theorems~5.5 and~5.9]{Schuch2010PEPS}.
    The assertion that every parent-kernel vector lies in the
    closure span remains separate.
\end{theorem}

\begin{proof}\leanok
    Regular $G$-injectivity makes every native closure vector
    nonzero. The local parent constraint therefore forces trivial
    plaquette holonomy. At the intersection of the two seams,
    this is precisely the commutation relation for the closure
    labels.

    A simply connected closed-cell region admits an integer lift.
    For commuting labels, the lift defines a vertex gauge which
    removes every internal closure operator. The remaining
    operators merely relabel the virtual boundary, so the original
    and inserted regional ground spaces agree. Every physical
    slice of the closed inserted contraction thus satisfies the
    original regional parent condition. The kernel of a positive
    sum is the intersection of the local kernels. Finally use
    linear independence of the commuting pair-conjugacy classes
    and their centralizer count.
\end{proof}
